
%\needspace{4\baselineskip}
\color{uniblue}{\subsection{ Устройство резервирования отказа выключателя (УРОВ)}}
\color{black}

\begin{enumerate}[label=\arabic{section}.\arabic{subsection}.\arabic{enumi}, labelsep=4pt, leftmargin=0pt, itemindent=57pt, itemsep=0pt, parsep=0pt] % если что можно крутить itemsep 0 или 5pt

\item
Устройство резервирования отказа выключателя (УРОВ) предназначено для выполнения функции ближнего резервирования защит. В случае отказа выключателя повреждённого присоединения УРОВ обеспечивает отключение смежных выключателей с целью селективного отделения повреждённого участка электрической сети.

\item
В случае присоединения линии к системе через два выключателя для каждого из них можно задействовать свою схему УРОВ. Далее приведено описание УРОВ для выключателя В1, логика УРОВ для второго выключателя аналогична.

\item
В устройстве применяется индивидуальный принцип реализации функции УРОВ. Пуск УРОВ происходит при срабатывании внутренних защит, выявлении неполнофазного режима, либо от внешних защит.

\item
Для выявления факта отказа выключателя предусмотрен контроль наличия тока через выключатель. Контроль осуществляется измерительным органом УРОВ, порог срабатывания которого задается уставкой \textbf{«Iср УРОВ В1»}.

\item
Для обеспечения надежного пуска УРОВ при кратковременном пропадании сигналов пуска до полного обесточивания отказавшего выключателя предусмотрен подхват от ИО УРОВ, ввод которого осуществляется с помощью программного ключа \textbf{«Фикс. УРОВ В1»}, сброс пуска произойдет только при возврате ИО УРОВ.

\item
Возможны два режима работы УРОВ:
\begin{enumerate}[label=\arabic*)]
    \item с автоматической проверкой исправности выключателя;
    \item с дублированным пуском.
\end{enumerate}

\item
В режиме «с автоматической проверкой исправности выключателя» программный ключ \textbf{«Контр. РПВ В1»} устанавливается в положение \textit{«Выведено»}, а программный ключ \textbf{«УРОВ В1 на себя»} -- \textit{«Введено»}. Данный режим выполняет функцию резервирования при обрыве цепи действия на отключение выключателя от защит. При появлении пуска УРОВ выдается команда на отключение «своего» выключателя, что позволяет исключить ложное срабатывание УРОВ на смежные выключатели. Выдержка времени действия 
УРОВ на «свой» выключатель регулируется уставкой \textbf{«Тср УРОВ В1 на себя»}. Если же отключения выключателя не происходит, значит неисправна не цепь от защит к ЭМО, а сам ЭМО. Тогда с проверкой срабатывания ИО УРОВ и с большей выдержкой времени \textbf{«Тср УРОВ В1»} формируется сигнал срабатывания УРОВ на смежные выключатели. Контроль наличия тока при действии УРОВ «на себя» может быть выведено программным ключом \textbf{«Контр. по I УРОВ В1 на себя»}.

\item
В режиме «с дублированным пуском» программный ключ \textbf{«Контр. РПВ В1»} устанавливается в положение \textit{«Введено»}, а программный ключ \textbf{«УРОВ В1 на себя»} – \textit{«Выведено»}. В данном режиме пуск УРОВ разрешается только при срабатывании защит с действием на отключение – пропадание входного сигнала ФСУ \textbf{«РПВ В1 от ЭМО»} является признаком срабатывания защит и работе цепи отключения. С выдержкой времени \textbf{«Тср УРОВ В1»} формируется сигнал срабатывания УРОВ на смежные выключатели.

\item
Входной сигнал ФСУ \textit{«Пуск УРОВ НН В1»} предусмотрен для пуска УРОВ от защит, срабатывание которых не сопровождается увеличением тока или сопровождается протеканием малых токов КЗ, недостаточных для срабатывания токовых ИО УРОВ. Например, от пускового органа УРОВ НН, предусматриваемого в защитах (авто)трансформатора, который срабатывает при КЗ на стороне НН (авто)трансформатора за токоограничивающим реактором. Данный пуск осуществляется с контролем положения выключателя \textit{«РПО В1 от ЭМВ»}.

\item
Информация о положении выключателя (сигналы «РПВ В1 от ЭМО» и «РПО В1 от ЭМВ») в логике УРОВ берется непосредственно от соответствующих электромагнитов управления.

\item
УРОВ В1 может быть оперативно выведено из работы от внешнего сигнала \textit{«Вывод УРОВ В1»}, назначенного на соответствующий дискретный вход или функциональную клавишу.

\item
Функциональная схема логики  УРОВ В1 приведена на рисунке \ref{pict_urov:UROV}. Функциональная схема логики УРОВ В2 выполняется аналогично схеме логики УРОВ В1.

\begin{figure}[H]
\centering
\includegraphics[width=0.9\textwidth,height=0.9\textheight,keepaspectratio]{\fbpath/761.3311 УРОВ/img/UROV.pdf}
\captionof{figure}{Функциональная схема логики УРОВ В1}\label{pict_urov:UROV}
\end{figure}

\item
Уставки УРОВ В1 представлены в таблице \ref{desc_urov:tbl1}. Уставки УРОВ В2 аналогичны уставкам УРОВ В1.

\small
\begin{longtable}{|>{\centering\arraybackslash}m{5.3cm}|>{\centering\arraybackslash}m{3.3cm}|>{\centering\arraybackslash}m{4.2cm}|>{\centering\arraybackslash}m{1.8cm}|>{\centering\arraybackslash}m{1cm}|}
\caption{Параметры для настройки УРОВ В1\hfill\vspace{-0.5\baselineskip}}\label{desc_urov:tbl1}\\ 
\hline
\rowcolor{gray!30}
Параметр & Условное обозначение на схеме & Значение/ Диапазон & Единица измерения & Шаг \\ 
\hline
\endfirsthead
\caption*{\hspace{3pt}\emph{Продолжение таблицы \ref{desc_urov:tbl1}\hfill\vspace{-0.5\baselineskip}}} \\ % сделано по ГОСТ 2.105 п.6.8.7
\hline
\rowcolor{gray!30}
Параметр & Условное обозначение на схеме & Значение/ Диапазон & Единица измерения & Шаг \\ 
%\hline
\endhead

%===m>УРОВ В1|УРОВ
\raggedright {'Description': ''} & \centering Rele_s_k_tom_vozvrata_max_float_1_x_sr & \centering 0,000 ... 40,000 & \centering  & \centering \arraybackslash 0,001 \\
\hline
\raggedright {'Description': ''} & \centering Pereklyuchatel_Vkh_2_bool_const_2_S & \centering Выведено \\ Введено & \centering --- & \centering \arraybackslash --- \\
\hline
\raggedright {'Description': ''} & \centering Nezavisimoe_RV_na_srabatyvanie_1_T_vremya & \centering 0 ... 100000 & \centering  & \centering \arraybackslash 1 \\
\hline
\raggedright {'Description': ''} & \centering Programmnyj_klyuch_3_s & \centering Введено \\ Выведено & \centering --- & \centering \arraybackslash --- \\
\hline
\raggedright {'Description': ''} & \centering Nezavisimoe_RV_na_srabatyvanie_2_T_vremya & \centering 0 ... 100000 & \centering  & \centering \arraybackslash 1 \\
\hline
\raggedright {'Description': ''} & \centering Programmnyj_klyuch_1_s & \centering Введено \\ Выведено & \centering --- & \centering \arraybackslash --- \\
\hline
\raggedright {'Description': ''} & \centering Programmnyj_klyuch_2_s & \centering Введено \\ Выведено & \centering --- & \centering \arraybackslash --- \\
\hline
\raggedright {'Description': ''} & \centering Pereklyuchatel_Vkh_2_bool_const_1_S & \centering Выведено \\ Введено & \centering --- & \centering \arraybackslash --- \\
\hline
%===m
\end{longtable}
\normalsize

\end{enumerate}